\documentclass{report}
\usepackage{amsfonts, amsmath}
\newcommand\R{{\mathbb R}}
\newcommand\Q{{\mathbb Q}}
\newcommand\N{{\mathbb N}}
\newcommand\Z{{\mathbb Z}}
\newcommand\Or{{\mathcal O}}
\newcommand\ve{\varepsilon}
%\pagestyle{empty}
\begin{document}
\large
\centerline{\Huge Fisica Matematica II}
\renewcommand{\thefootnote}{ } 
\renewcommand{\thefootnote}{\arabic{footnote}} 
\vskip.1cm
\centerline{\Large Soluzioni appello 17-06-2003}
\vskip1cm
\begin{enumerate}
\item Si controlla immediatamente che la condizione iniziale non \`e
una caratteristica. Ci si riconduce dunque al sistema
\[
\begin{split}
&\dot x=xz\\
&\dot y=1\\
&\dot z=1\\
&z(0)=\frac s2\\
&y(0)=s\\
&x(0)=1.
\end{split}
\]
Si ottengono facilmente i seguenti integrali del moto:
\[
z-y=c_1\;\quad \ln x-\frac {z^2}2=c_2.
\]
Il secondo si ottine integrando $\dot x=xz\dot z$, per separazione di
variabili.
Sostituendo le condizioni iniziali si ottine $c_1=-\frac s2$,
$c_2=\frac{s^2}8$. Dunque
\[
-\ln x+\frac {z^2}2=\frac {s^2}8=\frac 12(z-y)^2
\]
da cui segue
\[
z=\frac {2\ln x+y^2}{2y}.
\]
\item Si cerca una soluzione del tipo 
\[
u(x,t)=\sum_{k=1}^\infty c_k(t)\sin kx
\]
che soddisfa le condizioni al contorno. Sostituendo nell'equazione si
trova
\[
\dot c_k+k^2c_k=f_k\;\quad c_k(0)=h_k
\]
dove $f_k$ e $h_k$ sono i coefficienti della serie dei seni di
Fourier.
La soluzione \`e
\[
c_k(t)=(h_k-\frac {f_k}{k^2})e^{-k^2 t}+\frac{f_k}{k^2}.
\]
Si noti che se $f\in L^\infty$ allora gli $f_k$ sono limitati e la
serie con coefficienti $c_k$ converge uniformemente per ogni $t\geq
t_0>0$. Dunque si puo prendere il limite per $t\to\infty$ ottenendo
\[
v(x,t)=\sum_{k=1}^\infty\frac{f_k}{k^2}\sin kx.
\]
Chiaramente
\[
v_{xx}=-f(x)
\]
cio\`e $v$ \`e una soluzione stazionaria (non dipendente dal tempo)
della equazione originaria.
\item Se si introduce la relazione di equivalenza, in $\R^2$,
$(x,y)\sim(z,t)$ se e solo se
$(x,y)-(z,t)=nv+mw$ con $n,m\in\Z$ si ottine un toro
bidimensionale e la funzione $f$ \`e ben definita su di esso. In altre
parole $f$ \`e periodica sul reticolo formato dai vettori $v$, $w$
(che, per le persone prone alla geometria, pu\`o essere visto come il
ricoprimento universale del toro di cui sopra). Ad ogni modo ci\`o
significa che la funzione $f$ \`e limitata su $\R^2$ e dunque costante
per il teorema di Liuouville.
\item Prima di tutto si noti che ponendo $v(x,t)=u(x,t)-\sin x$, si ha
\[
\begin{split}
&v_{xx}=v_{tt}\\
&v(0,t)=0 \quad t>0\\
&v(x,0)=x-\sin x;\;v_t(x,0)=0 \quad x\geq 0.
\end{split}
\]
Secondo, si consideri il problema
\[
\begin{split}
&\tilde v_{xx}=\tilde v_{tt}\\
&\tilde v(x,0)=x-\sin x;\;\tilde v_t(x,0)=0 \quad x\in\R.
\end{split}
\]
Chiaramente ogni soluzione dispari di tale problema \`e una solu\-zione del
problema originario.
D'altro canto la soluzione dell'ulti\-mo problema \`e (formula di d'Alembert)
\[
\tilde v(x,t)=\frac 12\{ 2x-\sin (x+t)-\sin(x-t)\}=x-\sin x\cos t.
\]
che \`e chiaramente dispari.  Quindi
\[
u(x,t)=x+\sin x-\sin x\cos t.
\]
\end{enumerate}
\end{document}
